% Appendix F: Gesture Manifolds, Canonicalization, and Motor Equivalence

\section*{F.1 Introduction}

The previous appendices studied states. Human action, however, is not
fundamentally state-based — it is trajectory-based. A typed word is not
a sequence of letters; it is a path through a space of possibilities.
The central thesis is:
\[
  \text{Motor trajectory}
  \;\xrightarrow{C}\;
  \text{Canonical gesture}
  \;\xrightarrow{\Phi}\;
  \text{Symbol.}
\]

\begin{definition}{Gesture Path Space}{gesture-paths}
  A \emph{gesture} is a continuous path $g : [0,T] \to X$, where $X$
  encodes finger states, hand states, timing, pressure, and position.
  The \emph{gesture space} is $G = \mathrm{Path}([0,T], X)$.
\end{definition}

\begin{definition}{Gesture Manifold}{gest-mfld}
  A \emph{gesture manifold} is a smooth submanifold $M \subseteq G$.
  Skill corresponds to confinement within specialized gesture manifolds.
\end{definition}

\begin{definition}{Gesture Equivalence}{gest-equiv}
  $g_1 \sim g_2$ when they differ only by admissible perturbations
  (timing jitter, small spatial deformation, execution noise).
\end{definition}

\begin{definition}{Canonicalization}{canon-op}
  A \emph{canonicalization operator} is $C : G \to \hat{G}$ satisfying
  $C(g_1) = C(g_2)$ whenever $g_1 \sim g_2$.
  The \emph{canonical gesture space} is $\hat{G} = G/{\sim}$.
\end{definition}

\begin{definition}{Pivot}{pivot}
  A \emph{pivot} occurs at $t_p$ when
  $\dot{g}(t_p^-) \cdot \dot{g}(t_p^+) < 0$ (velocity reverses sign).
  Pivots correspond to syllabic peaks, typing transitions, musical accents.
\end{definition}

\begin{definition}{Gesture Energy}{gest-energy}
  $E(g) = \int_0^T \|\dot{g}(t)\|^2\,dt$.
  Motor learning tends toward lower-energy representatives within
  the same equivalence class.
\end{definition}

\begin{theorem}{Gesture Geodesics}{gest-geo}
  Among all equivalent motions, the geodesic (energy-minimizing)
  trajectory is most efficient. Repeated practice converges toward
  geodesic representatives.
\end{theorem}

\begin{definition}{Symbolic Projection}{sym-proj}
  A \emph{symbolic projection} is $\Phi : \hat{G} \to \Sigma$
  assigning a symbol to each canonical gesture class.
\end{definition}

\begin{theorem}{Gesture–Symbol Theorem}{gest-sym}
  Symbols are projections of gesture equivalence classes, not primitives.
\end{theorem}
\begin{proof}
  Canonicalization produces $[g]$; projection assigns $\Phi([g])$.
  Symbols arise from equivalence classes rather than individual trajectories.
\end{proof}

\begin{namedthm}{The Gesture Substrate Principle}
  A symbol is not a primitive thing. It is a projection of a motor
  equivalence class. Symbolic systems operate in the image of $\Phi$
  while discarding the richer fiber structure of $\hat{G}$.
\end{namedthm}
